\documentclass[11pt]{article}

% ---- packages ----
\usepackage[utf8]{inputenc}
\usepackage[T1]{fontenc}
\usepackage{lmodern}
\usepackage[margin=1in]{geometry}
\usepackage{amsmath,amssymb,amsthm,mathtools}
\usepackage{braket}
\usepackage{enumitem}
\usepackage{xcolor}
\usepackage[colorlinks=true,linkcolor=blue,urlcolor=blue]{hyperref}

% ---- macros (subset of the course preamble) ----
\newcommand{\R}{\mathbb{R}}
\newcommand{\C}{\mathbb{C}}
\newcommand{\E}{\mathbb{E}}
\newcommand{\Prob}{\mathbb{P}}
\newcommand{\ii}{\mathrm{i}}
\newcommand{\ee}{\mathrm{e}}
\DeclareMathOperator{\tr}{tr}
\DeclareMathOperator{\Tr}{Tr}
\DeclareMathOperator{\Var}{Var}
\DeclareMathOperator{\diag}{diag}
\newcommand{\norm}[1]{\left\lVert #1 \right\rVert}
\newcommand{\abs}[1]{\left\lvert #1 \right\rvert}
\newcommand{\bigo}{\mathcal{O}}
\newcommand{\id}{\mathbb{I}}
\newcommand{\dg}{^{\dagger}}
\newcommand{\proj}[1]{\ket{#1}\!\bra{#1}}
\newcommand{\CNOT}{\mathrm{CNOT}}
\newcommand{\QFT}{\mathrm{QFT}}
\newcommand{\defeq}{\vcentcolon=}

% ---- problem environment ----
\newlist{problems}{enumerate}{1}
\setlist[problems]{label=\textbf{Problem \arabic*.},leftmargin=*,itemsep=1.1em,topsep=0.6em}
\newcommand{\pts}[1]{\hfill{\normalfont\itshape(#1 pts)}}

\title{\textbf{Quantum Scientific Computing (QSC2026)}\\[0.3em]\large Homework~1}
\author{}
\date{Assigned: \qquad\qquad Due: \rule{3cm}{0.4pt}}

\begin{document}
\maketitle
\vspace{-2.2em}

\noindent\textbf{Instructions.} Eight problems drawn from Lectures~1, 3, and~4
(foundations; Grover search; the quantum Fourier transform and phase estimation).
Each is worth $10$ points. Show your
work: a correct number with no derivation earns little credit, and a clear
argument with a small arithmetic slip earns most of it. You may use any result
stated in the lecture notes, citing it by name. Collaboration is allowed on ideas,
but write up your solutions independently.

\bigskip
\noindent\textsc{Part I --- Foundations (Lecture 1)}
\begin{problems}

\item \textbf{Entanglement and no-cloning.}\pts{10}
Work entirely with pure states.
\begin{enumerate}[label=(\alph*),nosep,leftmargin=2em]
  \item Compute $\CNOT\ket{+}\ket{+}$ and decide whether the result is a product
        state or entangled, justifying your answer.
  \item Exhibit a unitary $U$ and a fixed ``blank'' state $\ket{s}$ for which
        $U\ket{0}\ket{s}=\ket{0}\ket{0}$ and $U\ket{1}\ket{s}=\ket{1}\ket{1}$, and
        explain why this does \emph{not} let you copy an unknown qubit
        $\alpha\ket0+\beta\ket1$---that is, why it is consistent with the no-cloning
        theorem.
\end{enumerate}

\item \textbf{Quantum teleportation.}\pts{10}
In the teleportation protocol, suppose Penelope's Bell measurement yields
$(m_1,m_2)=(1,0)$. Here the \emph{pre-correction state} means the state of
Odysseus's qubit immediately after Penelope's measurement and the transmission of
the two classical bits $(m_1,m_2)$, but \emph{before} Odysseus applies any
correction gate.
\begin{enumerate}[label=(\alph*),nosep,leftmargin=2em]
  \item Write down Odysseus's pre-correction state and verify it equals $Z\ket{\psi}$.
  \item State the correction Odysseus applies and confirm it restores $\ket{\psi}$.
\end{enumerate}

\item \textbf{The reversible two-bit adder.}\pts{10}
Using the half- and full-adder circuits of \S1.3.6 of the notes (``Worked example:
a reversible binary adder''):
\begin{enumerate}[label=(\alph*),nosep,leftmargin=2em]
  \item trace the two-bit adder on the input $1+2$ (that is, $A_1A_0=01$ and
        $B_1B_0=10$) and confirm the measured output is $011_2=3$;
  \item count the Toffoli and $\CNOT$ gates in the complete two-bit adder.
\end{enumerate}
\end{problems}

\bigskip
\noindent\textsc{Part II --- Query algorithms and Grover search (Lecture 3)}
\begin{problems}\setcounter{problemsi}{3}

\item \textbf{The four one-bit functions and Deutsch's algorithm.}\pts{10}
List all four functions $f:\{0,1\}\to\{0,1\}$ and classify each as constant or
balanced. For each, compute by hand the first-qubit measurement outcome produced
by the Deutsch algorithm, and confirm that the outcome distinguishes the constant
from the balanced functions in a single query.

\item \textbf{Iteration count and over-rotation.}\pts{10}
\begin{enumerate}[label=(\alph*),nosep,leftmargin=2em]
  \item For $N=2^{20}$, approximately how many iterations $k^{\ast}$ maximize the
        success probability? What happens qualitatively to the success probability
        if one instead runs $2k^{\ast}$ iterations? Explain using
        $p_k=\sin^2\!\bigl((2k+1)\theta/2\bigr)$.
  \item For $M=4$ marked items among $N=1024$, find $\sin(\theta/2)$, the optimal
        iteration count $k^{\ast}$, and compare it with the single-solution count
        for the same $N$.
\end{enumerate}

\item \textbf{Building the phase oracle.}\pts{10}
On three qubits, let the marked state be $x_0=\ket{101}$. Write the marking
reflection $R_{x_0}=\id-2\proj{x_0}$ as a doubly controlled $Z$ conjugated by $X$
gates, specifying which qubits carry the $X$ gates. Verify by hand that your
circuit multiplies $\ket{101}$ by $-1$ and leaves the other seven basis states
fixed (it suffices to check $\ket{101}$, one basis state that becomes $\ket{111}$
under the leading NOTs, and one that does not). Then run \emph{one} Grover iterate
on the uniform state $\ket{\psi_0}=H^{\otimes3}\ket{000}$ with this $x_0$ and report
the resulting amplitudes.
\end{problems}

\bigskip
\noindent\textsc{Part III --- Quantum Fourier transform and phase estimation (Lecture 4)}
\begin{problems}\setcounter{problemsi}{6}

\item \textbf{Exact phase estimation.}\pts{10}
Let $\varphi=0.011_2=\tfrac38$ and take $t=3$ ancilla qubits. Write the ancilla
state after the controlled powers in standard QPE, identify it as $\QFT\ket{k}$ for
the correct $k$, and conclude the measurement outcome. Why is the readout exact in
this case?

\item \textbf{How many ancillas?}\pts{10}
Using the precision/qubit-count relation of Lecture~4, determine how many ancilla
qubits $t$ are needed to estimate a phase $\varphi$ to $n=6$ bits of precision with
success probability at least $0.9$ (so $\epsilon=0.1$). State the general formula
you use and then evaluate it.
\end{problems}

\vfill
\noindent\footnotesize Problems adapted from the exercise sets of Lectures~1, 3,
and~4 of the Quantum Scientific Computing course notes.
\end{document}
